\subsection{Duality of extension modules and torsion products}

Let A be a ring, and let P and J be A-modules. For every complex C of A-modules, we constructed in A, X, p.~99, prop. 12 a canonical isomorphism of complexes

\[
\mu : \operatorname{Homgr} _ {\mathrm{A}} (\mathrm{C} \otimes_ {\mathrm{A}} \mathrm{P}, \mathrm{J}) \longrightarrow \operatorname{Homgr} _ {\mathrm{A}} (\mathrm{C}, \operatorname{Hom} _ {\mathrm{A}} (\mathrm{P}, \mathrm{J})) .
\]

Let M be a \(\Lambda\) module, and (C, p) a projective resolution of M. Consider the sequence of homomorphisms

\[
\begin{array}{r l} \operatorname{Ext} _ {\mathrm{A}} (\mathrm{M}, \operatorname{Hom} _ {\mathrm{A}} (\mathrm{P}, \mathrm{J})) & \xrightarrow {\varphi^ {- i}} \mathrm{H} (\operatorname{Homgr} _ {\mathrm{A}} (\mathrm{C}, \operatorname{Hom} _ {\mathrm{A}} (\mathrm{P}, \mathrm{J}))) \xrightarrow {\mathrm{H} (\mu) ^ {- 1}} \mathrm{H} (\operatorname{Homgr} _ {\mathrm{A}} (\mathrm{C} \otimes_ {\mathrm{A}} \mathrm{P}, \mathrm{J})) \\ & \xrightarrow {u} \operatorname{Homgr} _ {\mathrm{A}} (\mathrm{H} (\mathrm{C} \otimes_ {\mathrm{A}} \mathrm{P}), \mathrm{J}) \xrightarrow {v} \operatorname{Homgr} _ {\mathrm{A}} (\operatorname{Tor} ^ {\mathrm{A}} (\mathrm{M}, \mathrm{P}), \mathrm{J}), \end{array}
\]

where \(\varphi\) is the canonical isomorphism \(\varphi(\mathrm{C},\mathrm{Hom}_{\mathrm{A}}(\mathrm{P},\mathrm{J}))\) (A, X, p.~100, th. 1) \(u\) the canonical homomorphism \(\lambda(\mathrm{C}\otimes_{\Lambda}\mathrm{P},\mathrm{J})\) . (A, X, p.~82), and \(v\) is deduced from the canonical isomorphism \(\psi(\mathrm{C},\mathrm{P}):\operatorname{Tor}^{\Lambda}(\mathrm{M},\mathrm{P})\longrightarrow\mathrm{H}(\mathrm{C}\otimes_{\mathrm{A}}\mathrm{P})\) .

Let \((\mathbf{C}', p')\) another projective resolution of M. By A, X, p.~49, cor. of prop. 3, there exists a homotopy of complexes \(\alpha: \mathbf{C}' \to \mathbf{C}\) such that \(p \circ \alpha = p'\) . It follows from A, X, p.~103, prop. 2, that one has \(\mathrm{H}(\alpha \otimes 1_{\mathrm{P}}) \circ \psi(\mathrm{C}', \mathrm{P}) = \psi(\mathrm{C}, \mathrm{P})\) and \(\varphi(\mathrm{C}', \mathrm{R}) \circ \mathrm{H}(\mathrm{Homgr}(\alpha, 1_{\mathrm{R}})) = \varphi(\mathrm{C}, \mathrm{R})\) for every A-module R. From this it follows that the graded homomorphism of degree 0

\[
\theta (\mathrm{M}, \mathrm{P}): \operatorname{Ext} _ {\mathrm{A}} (\mathrm{M}, \operatorname{Hom} _ {\mathrm{A}} (\mathrm{P}, \mathrm{J})) \longrightarrow \operatorname{Homgr} _ {\mathrm{A}} \left(\operatorname{Tor} ^ {\mathrm{A}} (\mathrm{M}, \mathrm{P}), \mathrm{J}\right)
\]

the composite of the sequence of homomorphisms above is independent of the choice of the projective resolution (C,p) of M. By construction it is End \(_{A}\) (J)-linear.

The definition of the homomorphism \(\theta(\mathrm{M},\mathrm{P})\) is made explicit as follows. Let \(p\) an integer, \(\mathfrak{v}\) be an element of \(\operatorname{Ext}_{\Lambda}^{p}(\mathrm{M},\operatorname{Hom}_{\mathrm{A}}(\mathrm{P},\mathrm{J}))\) , \(\tau\) be an element of \(\operatorname{Tor}_{p}^{\mathrm{A}}(\mathrm{M},\mathrm{P})\) . Using the isomorphism \(\varphi(\mathrm{C},\operatorname{Hom}_{\mathrm{A}}(\mathrm{P},\mathrm{J}))\) , \(\mathfrak{v}\) is represented by a linear map \(u:\mathrm{C}_{p}\to\operatorname{Hom}_{\mathrm{A}}(\mathrm{P},\mathrm{J})\) such that \(u\circ d_{\mathrm{C}}=0\) ; similarly, using \(\psi(\mathrm{C},\mathrm{P})\) , \(\tau\) is represented by an element \(\sum c_{\mu}\otimes p_{\mu}\) of \(\mathrm{C}_{p}\otimes\mathrm{P}\) such that \(\sum d_{\mathrm{C}}(c_{\mu})\otimes p_{\mu}=0\) . One then has \(\theta(\mathrm{M},\mathrm{P})(\mathfrak{v})(\tau)=\sum u(c_{\mu})(p_{\mu})\) .

On the other hand, let \(\nu : C \otimes_A \operatorname{Hom}_A(P, J) \longrightarrow \operatorname{Homgr}_A(\operatorname{Homgr}_A(C, P), J)\) be the homomorphism that maps the element \(c \otimes h\) , for \(c \in C_p\) , \(h \in \operatorname{Hom}_A(P, J)\) on the homomorphism \(u \mapsto (-1)^p h(u(c))\) . It is graded of degree 0; it is bijective if each module \(C_p\) is free of finite type. One verifies without difficulty that it is a morphism of complexes.

Consider the sequence of homomorphisms

\[
\begin{array}{r l} \operatorname{Tor} ^ {\Lambda} (M, \operatorname{Hom} _ {\Lambda} (P, J)) & \xrightarrow {\psi} H (C \otimes_ {A} \operatorname{Hom} _ {\Lambda} (P, J)) \xrightarrow {H (\nu)} H (\operatorname{Homgr} _ {A} (\operatorname{Homgr} _ {A} (C, P), J)) \\ & \xrightarrow {w} \operatorname{Homgr} _ {\Lambda} (H (\operatorname{Homgr} _ {A} (C, P)), J) \xrightarrow {t} \operatorname{Homgr} _ {A} (\operatorname{Ext} _ {A} (M, P), J) \end{array}
\]

where \(\psi\) is the canonical isomorphism \(\psi(\mathrm{C},\mathrm{Hom}_{\mathrm{A}}(\mathrm{P},\mathrm{J}))\) , \(w\) the canonical homomorphism \(\lambda(\mathrm{Homgr}_{\mathrm{A}}(\mathrm{C},\mathrm{P}),\mathrm{J})\) (A, X, p.~82) and \(t\) is deduced from the canonical isomorphism \(\varphi(\mathrm{C},\mathrm{P})\) . One sees as above that the composed homomorphism

\[
\rho (\mathrm{M}, \mathrm{P}): \operatorname{Tor} ^ {\mathrm{A}} (\mathrm{M}, \operatorname{Hom} _ {\mathrm{A}} (\mathrm{P}, \mathrm{J})) \longrightarrow \operatorname{Homgr} _ {\mathrm{A}} \left(\operatorname{Ext} _ {\mathrm{A}} (\mathrm{M}, \mathrm{P}), \mathrm{J}\right)
\]

is independent of the choice of the resolution C; it is \(\operatorname{End}_{\mathrm{A}}(\mathrm{J})\) -linear. Let p be an integer, \(\xi \in \operatorname{Tor}_{p}^{\mathrm{A}}(\mathrm{M}, \operatorname{Hom}_{\mathrm{A}}(\mathrm{P}, \mathrm{J}))\) , \(\lambda \in \operatorname{Ext}_{\mathrm{A}}^{p}(\mathrm{M}, \mathrm{P}), \mathrm{J}\) ; if \(\xi\) is represented by means of \(\psi(\mathrm{C}, \operatorname{Hom}_{\Lambda}(\mathrm{P}, \mathrm{J}))\) by an element \(\sum c_{\mu} \otimes u_{\mu}\) of \(\mathrm{C} \otimes \operatorname{Hom}_{\Lambda}(\mathrm{P}, \mathrm{J})\) such that \(\sum d_{\mathrm{C}}(c_{\mu}) \otimes u_{\mu} = 0\) , and by \(\lambda\) by means of \(\varphi(\mathrm{C}, \mathrm{P})\) by a homomorphism \(\ell : C_{p} \to P\) such that \(\ell \circ d_{C} = 0\) , one has \(\rho(\mathrm{M}, \mathrm{P})(\xi)(\lambda) = (-1)^{p} \sum u_{\mu} (\ell(c_{\mu}))\) .

PROPOSITION 6. - Suppose the A-module J is injective; for every A-module N, set D(N) = \(\mathrm{Hom}_{\Lambda}(N,J)\) .

\begin{enumerate}
\def\labelenumi{\alph{enumi})}
\item
  Homomorphisms \(\theta^{i}(\mathrm{M},\mathrm{P}):\operatorname{Ext}_{\mathrm{A}}^{i}(\mathrm{M},\mathrm{D}(\mathrm{P}))\longrightarrow \mathrm{D}\big(\operatorname {Tor}_{i}^{\mathrm{A}}(\mathrm{M},\mathrm{P})\big)\) are bijective.
\item
  If the ring A is Noetherian and the A-module M is finitely generated, the homomorphisms \(\rho_{i}(\mathrm{M},\mathrm{P}):\operatorname{Tor}_{i}^{\mathrm{A}}(\mathrm{M},\mathrm{D}(\mathrm{P}))\longrightarrow \mathrm{D}\bigl (\operatorname {Ext}_{\mathrm{A}}^{i}(\mathrm{M},\mathrm{P})\bigr)\) are bijective.
\item
  By construction, the homomorphism \(\theta(\mathrm{M},\mathrm{P})\) is bijective as soon as \(\lambda(\mathrm{C}\otimes_{\mathrm{A}}\mathrm{P},\mathrm{J})\) is bijective, which is the case when J is injective (A, X, p.~85, cor. 2).
\item
  Choose the resolution C so that each module \(C_{p}\) is free of finite type (A, X, p.~53, prop. 6). Then the homomorphism v is bijective, and the same is true of \(\lambda(\text{Homgr}_{\text{A}}(\text{C}, \text{P}), \text{J})\), since J is injective; therefore \(\rho(\text{M}, \text{P})\) is bijective.
\end{enumerate}

Remarks.--- 1) For every homomorphism \(f: \mathrm{N} \to \mathrm{N}'\) of A-modules, denote \(\mathrm{D}(f): \mathrm{D}(\mathrm{N}') \to \mathrm{D}(\mathrm{N})\) the homomorphism \(\mathrm{Hom}(f, \mathbf{1}_{\mathrm{J}})\) . Let \(u: \mathrm{M} \to \mathrm{M}'\) and \(v: \mathrm{P} \to \mathrm{P}'\) homomorphisms of A-modules. Choose projective resolutions (C, p) of M and (C', p') of M', and a morphism of complexes \(\tilde{u}: \mathrm{C} \to \mathrm{C}'\) such that \(p' \circ \tilde{u} = u \circ p\) (A, X, p.~49, prop. 3). The diagram

\[
\begin{array}{c c c} \operatorname{Homgr} _ {\mathrm{A}} (\mathrm{C} ^ {\prime} \otimes_ {\mathrm{A}} \mathrm{P} ^ {\prime}, \mathrm{J}) & \xrightarrow {\mu^ {\prime}} & \operatorname{Homgr} _ {\mathrm{A}} (\mathrm{C} ^ {\prime}, \operatorname{Hom} _ {\mathrm{A}} (\mathrm{P} ^ {\prime}, \mathrm{J})) \\ \operatorname{Hom} (\tilde {u} \otimes v, 1 _ {\mathrm{J}}) \Bigg \downarrow & & \Bigg \downarrow \operatorname{Hom} (\tilde {u}, \operatorname{Hom} (v, 1)) \\ \operatorname{Homgr} _ {\mathrm{A}} (\mathrm{C} \otimes_ {\mathrm{A}} \mathrm{P}, \mathrm{J}) & \xrightarrow {\mu} & \operatorname{Homgr} _ {\mathrm{A}} (\mathrm{C}, \operatorname{Hom} _ {\mathrm{A}} (\mathrm{P}, \mathrm{J})) \end{array}
\]

where \(\mu\) and \(\mu'\) are the canonical homomorphisms, is commutative; one then deduces from A. X, p.~103, prop. 2 a commutative diagram

\[
\begin{array}{c c c} \operatorname{Ext} _ {\Lambda} ^ {i} (\mathrm{M} ^ {\prime}, \mathrm{D} (\mathrm{P} ^ {\prime})) & \xrightarrow {\theta^ {i} (\mathrm{M} ^ {\prime} , \mathrm{P} ^ {\prime})} & \mathrm{D} (\operatorname{Tor} _ {i} ^ {\mathrm{A}} (\mathrm{M} ^ {\prime}, \mathrm{P} ^ {\prime})) \\ \operatorname{Ext} ^ {i} (u, \mathrm{D} (v)) \Bigg \downarrow & & \Bigg \downarrow \mathrm{D} (\operatorname{Tor} _ {i} (u, v)) \\ \operatorname{Ext} _ {\Lambda} ^ {i} (\mathrm{M}, \mathrm{D} (\mathrm{P})) & \xrightarrow {\theta^ {i} (\mathrm{M} , \mathrm{P})} & \mathrm{D} (\operatorname{Tor} _ {i} ^ {\Lambda} (\mathrm{M}, \mathrm{P})) \end{array}
\]

Let \(w: \mathrm{P}'' \to \mathrm{P}\) a homomorphism of A-modules; in an analogous manner we obtain a commutative diagram

\[
\begin{array}{c c c} \operatorname{Tor} _ {i} ^ {\mathrm{A}} (\mathrm{M}, \mathrm{D} (\mathrm{P})) & \xrightarrow {\rho_ {i} (\mathrm{M} , \mathrm{P})} & \mathrm{D} (\operatorname{Ext} _ {\mathrm{A}} ^ {i} (\mathrm{M}, \mathrm{P})) \\ \operatorname{Tor} _ {i} (u, \mathrm{D} (w)) \Bigg \downarrow & & \Bigg \downarrow \mathrm{D} (\operatorname{Ext} ^ {i} (u, w)) \\ \operatorname{Tor} _ {i} ^ {\mathrm{A}} (\mathrm{M} ^ {\prime}, \mathrm{D} (\mathrm{P} ^ {\prime \prime})) & \xrightarrow {\rho_ {i} (\mathrm{M} ^ {\prime} , \mathrm{P} ^ {\prime \prime})} & \mathrm{D} (\operatorname{Ext} _ {\mathrm{A}} ^ {i} (\mathrm{M} ^ {\prime}, \mathrm{P} ^ {\prime \prime})) \end{array} .
\]

\begin{enumerate}
\def\labelenumi{\arabic{enumi})}
\setcounter{enumi}{1}
\tightlist
\item
  Let
\end{enumerate}

\[
0 \to \mathrm{M} ^ {\prime} \xrightarrow {j} \mathrm{M} \xrightarrow {q} \mathrm{M} ^ {\prime \prime} \to 0\tag{ℰ}
\]

an exact sequence of A-modules. The homomorphism \(\mathrm{L}(q):\mathrm{L}(\mathrm{M})\to\mathrm{L}(\mathrm{M}^{\prime\prime})\) induced on the canonical free resolutions is surjective, and the complex \(\operatorname{Ker}\mathrm{L}(q)\) defines a projective resolution of \(M^{\prime}\) Applying prop. 3 of A, X, p.~104 to the exact sequence \(0\to\operatorname{Ker}\mathrm{L}(q)\to\mathrm{L}(\mathrm{M})\to\mathrm{L}(\mathrm{M}^{\prime\prime})\to0\) , we obtain commutative diagrams

\[
\begin{array}{c c c} \operatorname{Ext} _ {\mathrm{A}} ^ {i} (\mathrm{M} ^ {\prime}, \mathrm{D} (\mathrm{P})) & \xrightarrow {\theta^ {i} (\mathrm{M} ^ {\prime} , \mathrm{P})} & \mathrm{D} (\operatorname{Tor} _ {i} ^ {\mathrm{A}} (\mathrm{M} ^ {\prime}, \mathrm{P})) \\ \delta^ {i} (\mathcal {E}, \mathrm{D} (\mathrm{P})) \Bigg \downarrow & & \Bigg \downarrow (- 1) ^ {i + 1} \mathrm{D} (\partial_ {i + 1} (\mathcal {E}, \mathrm{P})) \\ \operatorname{Ext} _ {\mathrm{A}} ^ {i + 1} (\mathrm{M} ^ {\prime \prime}, \mathrm{D} (\mathrm{P})) & \xrightarrow {\theta^ {i + 1} (\mathrm{M} ^ {\prime \prime} , \mathrm{P})} & \mathrm{D} (\operatorname{Tor} _ {i + 1} ^ {\mathrm{A}} (\mathrm{M} ^ {\prime \prime}, \mathrm{P})) \end{array}
\]

\[
\begin{array}{c c c} \operatorname{Tor} _ {i + 1} ^ {\mathrm{A}} (\mathrm{M} ^ {\prime \prime}, \mathrm{D} (\mathrm{P})) & \xrightarrow {\rho_ {i + 1} (\mathrm{M} ^ {\prime \prime} , \mathrm{P})} & \mathrm{D} (\operatorname{Ext} _ {\mathrm{A}} ^ {i + 1} (\mathrm{M} ^ {\prime \prime}, \mathrm{P})) \\ \partial_ {i + 1} (\mathcal {E}, \mathrm{D} (\mathrm{P})) \Bigg \downarrow & & \Bigg \downarrow (- 1) ^ {i + 1} \mathrm{D} (\delta^ {i} (\mathcal {E}, \mathrm{P})) \\ \operatorname{Tor} _ {i} ^ {\mathrm{A}} (\mathrm{M} ^ {\prime}, \mathrm{D} (\mathrm{P})) & \xrightarrow {\rho_ {i} (\mathrm{M} ^ {\prime} , \mathrm{P})} & \mathrm{D} (\operatorname{Ext} _ {\mathrm{A}} ^ {i} (\mathrm{M} ^ {\prime}, \mathrm{P})) \end{array}
\]

Let

\[
0 \to \mathrm{P} ^ {\prime} \to \mathrm{P} \to \mathrm{P} ^ {\prime \prime} \to 0\tag{F}
\]

an exact sequence of A-modules; since the A-module J is injective, one deduces an exact sequence

\[
0 \to \mathrm{D} (\mathrm{P} ^ {\prime \prime}) \to \mathrm{D} (\mathrm{P}) \to \mathrm{D} (\mathrm{P} ^ {\prime}) \to 0.\tag{\(\left(\mathcal{D}(\mathcal{F})\right)\)}
\]

Applying A, X, p.~104, prop. 3 and p.~106, prop. 4 to the exact sequences (F) and (D(F)), one obtains in an analogous manner commutative diagrams.

\[
\begin{array}{c c c} \operatorname{Ext} _ {\Lambda} ^ {i} (\mathrm{M}, \mathrm{D} (\mathrm{P} ^ {\prime})) & \xrightarrow {\theta^ {i} (\mathrm{M} , \mathrm{P} ^ {\prime})} & \mathrm{D} (\operatorname{Tor} _ {i} ^ {\mathrm{A}} (\mathrm{M}, \mathrm{P} ^ {\prime})) \\ \delta^ {i} (\mathrm{M}, \mathrm{D} (\mathcal {F})) \Bigg \downarrow & & \Bigg \downarrow (- 1) ^ {i + 1} \mathrm{D} (\partial_ {i + 1} (\mathrm{M}, \mathcal {S})) \\ \operatorname{Ext} _ {\mathrm{A}} ^ {i + 1} (\mathrm{M}, \mathrm{D} (\mathrm{P} ^ {\prime \prime})) & \xrightarrow {\theta^ {i + 1} (\mathrm{M} , \mathrm{P} ^ {\prime \prime})} & \mathrm{D} (\operatorname{Tor} _ {i + 1} ^ {\mathrm{A}} (\mathrm{M}, \mathrm{P} ^ {\prime \prime})) \end{array}
\]

\[
\begin{array}{c c c} \operatorname{Tor} _ {i + 1} ^ {\mathrm{A}} (\mathrm{M}, \mathrm{D} (\mathrm{P} ^ {\prime})) & \xrightarrow {\rho_ {i + 1} (\mathrm{M} , \mathrm{P} ^ {\prime})} & D (\operatorname{Ext} _ {\Lambda} ^ {i + 1} (\mathrm{M}, \mathrm{P} ^ {\prime})) \\ \partial_ {i + 1} (\mathrm{M}, \mathrm{D} (\mathcal {F})) \Bigg \downarrow & & \Bigg \downarrow (\sim 1) ^ {i} \mathrm{D} (\delta^ {i} (\mathrm{M}, \mathcal {F})) \\ \operatorname{Tor} _ {i} ^ {\mathrm{A}} (\mathrm{M}, \mathrm{P} ^ {\prime \prime}) & \xrightarrow {\rho_ {i} (\mathrm{M} , \mathrm{P} ^ {\prime \prime})} & D (\operatorname{Ext} _ {\mathrm{A}} ^ {i} (\mathrm{M}, \mathrm{P} ^ {\prime \prime})) \end{array}
\]

